\documentclass{article}
\usepackage[zh]{homework}
\usepackage{needspace}
% 题目 1--10, 原作业说明及课外阅读忠实转录自 HW7.pdf; 以下承诺要求为原文件内容的引用.

\config{Algebra-1}{2026 秋季}{7}

\begin{document}

请在提交的作业正文最前面写下以下承诺:
\begin{quote}
  我承诺, 本次作业没有使用大语言模型.
\end{quote}

截止时间: \textbf{2026-11-13} ({\heiti 周五}) \textbf{23:59}

\paragraph{约定} $\mathbf{Ab}$ 表示阿贝尔群及群同态组成的范畴. 环均有乘法单位元, 环同态保持单位元. 复合 $g\circ f$ 表示先做 $f$, 再做 $g$.

本次作业共 10 题, 第 1--9 题必做, 第 10 题选做. 可使用课堂已证明的结果及前面习题的结论, 但不得直接引用当前小问要求证明的结论. 上述承诺适用于所提交的习题解答. 文末课外阅读无需提交, 可以使用大语言模型 (LLM) 辅助阅读.

\section*{习题}

\begin{problem}
设 $A$ 为阿贝尔群. 若存在有限子集 $S\subseteq A$ 使得 $A=\langle S\rangle$, 则称 $A$ 是有限生成的. 判断下列三个加法群是否有限生成. 若是, 给出一组有限生成元; 若否, 证明任取有限多个元素都不能生成全群.
\begin{enumerate}
  \item[(a)] $\Z\oplus(\Z/6\Z)$, 加法按坐标定义.
  \item[(b)] $(\Q,+)$, 即有理数加法群.
  \item[(c)] $\Z^{(\N)}$, 即只有有限个非零项的整数序列组成的群, 加法逐项定义.
\end{enumerate}
\end{problem}

\begin{problem}
设 $\mathcal C$ 是预加性范畴, 即每个态射集都是阿贝尔群, 且复合对两边的加法都满足分配律. 记 $0_{X,Y}$ 为 $\Hom_{\mathcal C}(X,Y)$ 的零元.
\begin{enumerate}
  \item[(a)] 证明以下三条等价: $X$ 是初对象; $X$ 是终对象; $\id_X=0_{X,X}$. 若 $\mathcal C$ 有零对象 $O$, 证明 $X\to O\to Y$ 的复合等于 $0_{X,Y}$.
  \item[(b)] 取 $\mathbf{Ab}$ 中由所有非零阿贝尔群构成的满子范畴. 说明它为何仍是预加性范畴, 并用第 (a) 问判断它是否有初对象或终对象.
\end{enumerate}
\end{problem}

\begin{problem}
设 $S$ 为任意集合. 阿贝尔群 $F$ 连同集合映射 $\iota:S\to F$, 若满足以下泛性质, 就称为 $S$ 上的自由阿贝尔群: 对任意阿贝尔群 $B$ 和集合映射 $u:S\to B$, 存在唯一群同态 $h:F\to B$ 使 $h\circ\iota=u$. 本题逐步构造这样的 $(F,\iota)$.
\begin{enumerate}
  \item[(a)] 对函数 $f:S\to\Z$, 定义其支撑为 $\operatorname{supp}(f)=\{s\in S:f(s)\ne0\}$, 并令
  \[
    \Z^{(S)}=\{f:S\to\Z:\operatorname{supp}(f)\text{ 是有限集}\}.
  \]
  在 $\Z^{(S)}$ 上定义逐点加法 $(f+g)(s)=f(s)+g(s)$. 证明 $\operatorname{supp}(f+g)\subseteq\operatorname{supp}(f)\cup\operatorname{supp}(g)$, 写出零元与加法逆元, 并验证 $\Z^{(S)}$ 是阿贝尔群.
  \item[(b)] 对每个 $s\in S$, 令
  \[
    e_s(t)=\begin{cases}
      1,&t=s,\\
      0,&t\ne s,
    \end{cases}
    \qquad \iota:S\to\Z^{(S)},\quad s\mapsto e_s.
  \]
  验证 $e_s\in\Z^{(S)}$, 并证明每个 $f\in\Z^{(S)}$ 都满足
  \[
    f=\sum_{s\in\operatorname{supp}(f)}f(s)e_s.
  \]
  再证明系数的唯一性: 若 $f=\sum_{s\in S}n_se_s$, 其中 $n_s\in\Z$ 且只有有限多个 $n_s$ 非零, 则 $n_s=f(s)$ 对所有 $s\in S$ 成立.
  \item[(c)] 设 $B$ 为阿贝尔群, $h,g:\Z^{(S)}\to B$ 为群同态. 利用第 (b) 问和同态保持整数倍的性质, 证明若 $h\circ\iota=g\circ\iota$, 则 $h=g$. 由此说明, 满足上述泛性质中条件的同态至多有一个.
  \item[(d)] 给定集合映射 $u:S\to B$, 定义
  \[
    h(f)=\sum_{s\in\operatorname{supp}(f)}f(s)u(s).
  \]
  说明这个公式有意义. 对 $f,g\in\Z^{(S)}$, 在共同的有限集 $T=\operatorname{supp}(f)\cup\operatorname{supp}(g)$ 上改写求和, 证明 $h(f+g)=h(f)+h(g)$. 再验证 $h\circ\iota=u$, 从而完成泛性质的证明. 最后说明, 映射
  \[
    \Hom_{\mathbf{Ab}}(\Z^{(S)},B)\longrightarrow\Hom_{\mathbf{Set}}(S,B),\qquad h\longmapsto h\circ\iota
  \]
  是双射, 其中右侧表示从 $S$ 到 $B$ 的底层集合的所有映射.
\end{enumerate}
\end{problem}

\begin{problem}
沿用第 3 题的记号. 给定集合映射 $u:S\to T$, 由自由阿贝尔群的泛性质, 存在唯一群同态 $u_*:\Z^{(S)}\to\Z^{(T)}$ 将 $e_s$ 送到 $e_{u(s)}$.
\begin{enumerate}
  \item[(a)] 用 $a\in\Z^{(S)}$ 在 $u^{-1}(t)$ 上的取值写出 $(u_*a)(t)$. 说明当纤维 $u^{-1}(t)$ 无限时公式为何仍有意义. 对 $v:T\to U$, 证明 $(v\circ u)_*=v_*\circ u_*$.
  \item[(b)] 取 $S=\{1,2,3\}$, $T=\{x,y\}$, 且 $u(1)=u(2)=x$, $u(3)=y$. 求 $u_*$ 的像及其核的一组基.
\end{enumerate}
\end{problem}

\begin{problem}
设 $\pi:\Z^2\to\Z/12\Z$ 为群同态, 它将 $(1,0)$ 送到 $[4]$, $(0,1)$ 送到 $[6]$.
\begin{enumerate}
  \item[(a)] 写出 $\pi(a,b)$, 列出像中的全部元素, 并求一个映到 $[2]$ 的整数对.
  \item[(b)] 求 $\ker\pi$ 的一组基, 并据此写出映到 $[2]$ 的所有整数对. 说明 $\Z^2/\ker\pi$ 是几阶循环群.
\end{enumerate}
\end{problem}

\Needspace{8\baselineskip}
\begin{problem}
设 $F(S)$ 是集合 $S$ 上的自由群, $F(S)^{\mathrm{ab}}=F(S)/[F(S),F(S)]$, 这里 $[F(S),F(S)]$ 是交换子生成的子群. 记 $\bar s$ 为生成元 $s$ 在商群中的像.
\begin{enumerate}
  \item[(a)] 利用自由群和阿贝尔化的泛性质, 构造同态 $F(S)^{\mathrm{ab}}\to\Z^{(S)}$, 使每个 $\bar s$ 映到 $e_s$; 再利用第 3 题的泛性质构造反向同态.
  \item[(b)] 检查两个复合在生成元上的作用, 证明这两个同态互逆. 当 $S$ 为无限集时, 举出 $\prod_{s\in S}\Z$ 中不属于 $\Z^{(S)}$ 的一个元素.
\end{enumerate}
\end{problem}

\begin{problem}
令 $C_n=\Z/n\Z$ ($n\ge2$), 在 $\End_{\mathbf{Ab}}(C_n)$ 中以逐点加法为加法, 以复合为乘法.
\begin{enumerate}
  \item[(a)] 证明 $f\mapsto f([1])$ 给出 $\End_{\mathbf{Ab}}(C_n)\to\Z/n\Z$ 的环同构. 写出逆映射, 并验证复合怎样对应剩余类乘法.
  \item[(b)] 在 $C_{15}$ 中, 分别求由 $f([1])=[6]$ 和 $g([1])=[4]$ 确定的自同态的核与像, 判断它们是否为自同构. 对可逆者写出逆映射.
\end{enumerate}
\end{problem}

\begin{problem}
对整数 $n\ge2$, 定义 $\varphi(n)=\#\{a\in\Z:1\le a\le n,\ \gcd(a,n)=1\}$.
\begin{enumerate}
  \item[(a)] 用 Bézout 等式证明 $[a]\in(\Z/n\Z)^\times$ 当且仅当 $\gcd(a,n)=1$. 设 $p_1,\ldots,p_r$ 是 $n$ 的全部不同素因子, 用计数证明
  \[
    \varphi(n)=n\prod_{i=1}^r\left(1-\frac1{p_i}\right).
  \]
  提示: 在 $\{1,\ldots,n\}$ 中数出每个 $p_i$ 的倍数. 对于同时是几个 $p_i$ 的倍数的整数, 计数时如何处理?
  \item[(b)] 设 $n\ge2$, $a\in\Z$ 且 $\gcd(a,n)=1$. 对单位群 $(\Z/n\Z)^\times$ 应用 Lagrange 定理, 证明欧拉定理
  \[
    a^{\varphi(n)}\equiv1\pmod n.
  \]
  当 $n=p$ 为素数时, 由此推出费马小定理: 若 $p\nmid a$, 则 $a^{p-1}\equiv1\pmod p$. 进一步证明, 对任意整数 $a$ 都有 $a^p\equiv a\pmod p$, 并说明这一表述与前一表述等价.
\end{enumerate}
\end{problem}

\begin{problem}
若 $x^k=0$ 对某个整数 $k\ge1$ 成立, 则称 $x$ 为幂零元; 若 $e^2=e$, 则称 $e$ 为幂等元.
\begin{enumerate}
  \item[(a)] 求 $\Z/12\Z$ 中的全部幂零元和全部幂等元, 并说明为什么没有遗漏.

  提示: 可分别在模 $3$ 和模 $4$ 下考察 $x^k=0$ 与 $x(x-1)=0$.
  \item[(b)] 证明整环的幂等元只有 $0$ 和 $1$. 从第 (a) 问选一个 $e\ne0,1$ 的幂等元, 用它写出两个非零剩余类, 其乘积为零.
\end{enumerate}
\end{problem}

\begin{problem}
({\heiti 选做}) 给定含幺环 $R$, 构造只有一个对象 $\ast$ 的范畴 $\mathcal C_R$, 其态射集是 $R$, 规定 $a\circ b=ab$.
\begin{enumerate}
  \item[(a)] 验证 $\mathcal C_R$ 是预加性范畴, 指出恒等态射. 反过来, 从单对象预加性范畴的态射群与复合运算恢复一个含幺环 (这里允许零环).
  \item[(b)] 预加性范畴之间的函子若在每个态射群上诱导加法群同态, 就称为加性函子. 证明从 $\mathcal C_R$ 到 $\mathcal C_T$ 的加性函子与含幺环同态 $R\to T$ 一一对应, 并分别指出函子的哪条性质给出了保乘法和保单位元.
\end{enumerate}
\end{problem}

\section*{课外阅读: 有限除环都是域}

{\heiti 本节无需提交, 可以使用大语言模型} (\textbf{LLM}) {\heiti 辅助阅读.} 课堂上证明了有限整环中每个非零元可逆. Wedderburn 小定理从另一个方向提问: 若有限环 $D$ 中每个非零元都可逆, 为什么乘法一定交换? 这里假设 $1\ne0$.

阅读 Reed College 讲义 \href{https://people.reed.edu/~jerry/332/21cyclo.pdf\#page=5}{\textit{Cyclotomic Polynomials}, 第 4 节, 定理 4.1, 第 5--6 页}. 先找出群 $D^\times$ 在 $D$ 上的共轭作用, 再认出类方程中的两类项: 中心里的单点轨道, 以及非中心元素的轨道. 分圆多项式 $\Phi_n$ 的定义可回看第 2 节.

\end{document}
